\documentclass[11pt]{article}
\usepackage{icmlko}
\begin{document}

\title{집 안에서는 아무 일도 없었다 --- 축 하나가 집 밖 전이의 절반을 먹는데
계기 열하나가 그걸 못 본다}
\author{월드모델 실험실}
\date{노트 441 · 2026년 8월 2일}
\maketitle

\begin{abstract}
노트 440 이 남긴 물음은 ``KR 과 CN 을 무엇이 가르나'' 였다. 레코드를 열어
보니 후보가 둘이었고 --- \textbf{CN 은 \texttt{entry\_friction} 이 통째로
없다}(KR 은 100\%) --- 그 축 하나만 가려 봤다. 사전 등록은 편지대로
통과했다. \textbf{그런데 예측한 이유가 아니었다.} grp 이득이 준 것은 grp
팔이 내려서가 아니라 \textbf{무표지 팔이 $+0.2330$ 에서 $+0.5504$ 로 뛰어서}
였다. 축 하나가 KR 전이를 \emph{망치고} 있었고 grp 의 이득은 그 손해를
되돌리던 것이다. 대조(\texttt{venue\_prominence})는 반대로 두 팔을 다 내렸다.

기계는 뚜렷했다. 학습 안에서 축과 라벨의 부호를 다 재 보니 \textbf{다섯 중
\texttt{entry\_friction} 만 부호가 안 맞는다} --- 모바일 $+0.648$ 대 만화
$-0.158$ · KR $-0.320$, 양수 3/7. 그래서 집 안에서도 같은 게 보여야 한다고
등록했다. \textbf{반증됐다.} LODO 열하나에서 부호와 효과의 상관은 예측한
음수가 아니라 \textbf{rho $+0.714$} 였고, 무엇보다 \textbf{집안 효과는 전부
잡음 크기}였다(최대 $|0.0568|$).

그래서 집 밖 효과는 집안 계기가 아니라 \textbf{다른 집 밖 짝}으로
검증했다. 비게임 앱은 짝 안 부호가 $-0.566$ 으로 KR 보다 더 극단이다.
\textbf{복제됐고 더 컸다}: $-0.1389 \to +0.2799$, \textbf{$+0.4188$}.
대조(\texttt{goods\_scale})는 다시 반대로 갔다($-0.1148$).

\begin{center}
\begin{tabular}{lrrr}
\hline
 & 무표지 & grp & 이득 \\
\hline
KR 만화 · 기준 & $+0.2330$ & $+0.5015$ & $+0.2685$ \\
\textbf{KR 만화 $-$ entry\_friction} & $\mathbf{+0.5504}$ & $+0.5241$ & $-0.0263$ \\
비게임 앱 · 기준 & $-0.1389$ & $+0.1047$ & $+0.2436$ \\
\textbf{비게임 앱 $-$ entry\_friction} & $\mathbf{+0.2799}$ & $+0.2394$ & $-0.0405$ \\
\hline
\end{tabular}
\end{center}

\textbf{두 짝 모두에서 축을 빼면 grp 이득이 사라진다.} grp 이 집 밖에서
산 것은 새 정보가 아니라 \emph{이 축 하나의 되돌림}이었다.
\end{abstract}

\begin{figure}[t]
\centering
\includegraphics[width=\textwidth]{figs/f1.pdf}
\caption{\textbf{가 · 나} --- 집 밖 두 짝에서 \texttt{entry\_friction} 을
예보 때만 가린 결과. 두 짝 다 \emph{무표지} 팔이 뛰고 grp 이득은 $0$ 아래로
내려간다 --- grp 이 집 밖에서 산 것은 새 정보가 아니라 이 축 하나의
되돌림이었다. 대조로 뺀 축(\texttt{venue\_prominence} · \texttt{goods\_scale})은
두 짝 다 반대로 간다. \textbf{다} --- 같은 가림을 집 안(LODO) 열하나에서
잰 것. 사전 등록한 기울기는 음수였는데 실제는 rho $+0.714$ 이고, 무엇보다
\emph{전부 잡음 띠 안팎의 작은 값}이다. 위의 점선 둘이 집 밖 두 짝의
자리 --- 집안 최대치의 $5.6$배 · $7.4$배다.}
\label{fig:one}
\end{figure}

\section{편지로는 통과하고 이유로는 틀렸다}

등록한 문장은 ``KR 에서 \texttt{entry\_friction} 하나만 가리면 grp 이득이
$+0.15$ 아래로 떨어진다'' 였다. 결과는 $-0.0263$ --- 통과다. 하지만 등록할
때 머리에 있던 그림은 \emph{그 축이 전이를 실어 나르고 있어서 빼면 grp 팔이
같이 내려앉는다} 였다. 실제로는 grp 팔이 $+0.5015 \to +0.5241$ 로 거의 안
움직였다. 이득이 준 것은 \textbf{바닥이 올라와서}다.

노트 133 이 금지한 것이 정확히 이 자리다. 숫자가 등록한 문턱을 넘었다고
해서 등록할 때 상상한 기계가 맞은 것이 아니다.

\textbf{대조가 일반화를 막는다.} 아무 축이나 빼서 되는 일이라면 이야기가
아니다. \texttt{venue\_prominence} 를 대신 가리니 무표지가 $+0.2330 \to
+0.0755$ 로 \emph{내려갔고} 이득은 $+0.3162$ 로 \emph{올라갔다}. 축마다
방향이 다르다.

\section{기계 --- 다섯 중 하나만 눈금이 뒤집힌다}

학습 도메인 안에서 축과 라벨의 스피어만을 전부 재 봤다.

\begin{center}
\begin{tabular}{lrrr}
\hline
축 & 학습 평균 & 양수 & KR \\
\hline
\texttt{target\_breadth} & $+0.292$ & 9/11 & $+0.471$ \\
\texttt{venue\_prominence} & $+0.128$ & 9/10 & $+0.162$ \\
\texttt{goods\_scale} & $+0.202$ & 6/8 & $+0.389$ \\
\texttt{media\_push} & $+0.247$ & 5/6 & --- \\
\textbf{\texttt{entry\_friction}} & $+0.074$ & \textbf{3/7} & $\mathbf{-0.320}$ \\
\hline
\end{tabular}
\end{center}

넷은 도메인이 바뀌어도 같은 쪽을 가리킨다. \texttt{entry\_friction} 만
모바일에서 $+0.648$, 만화에서 $-0.158$, 애니에서 $-0.124$, KR 에서
$-0.320$ 이다. 폭이 $0.97$ 이다. 모델은 큰 양수 쪽을 배워 KR 에 \emph{거꾸로}
쓴다 --- 이것이 노트 419 가 ``grp 은 도메인마다 눈금이 달라 거친 도메인
표지로 일한다'' 고 했던 것의 이름 붙은 범인이다.

\section{집 안에서 반증됐다}

기계가 맞다면 집 안에서도 보여야 한다. 등록했다: \emph{축을 빼면 도메인 안
부호가 음수인 곳은 오르고 양수인 곳은 내린다 --- 상관이 음수라야 한다.}

\textbf{반증됐다.} LODO 열하나에서 rho $\mathbf{+0.714}$ (p $0.071$ · n 7)
--- 부호가 정확히 거꾸로다. 모바일($+0.648$)이 $+0.0202$ 로 \emph{오르고}
만화($-0.158$)가 $-0.0440$ 으로 \emph{내렸다}.

\textbf{그리고 더 중요한 것.} 집안 효과는 전부 잡음 크기다. 제일 큰 것이
팝업 $-0.0568$ 인데 \textbf{KR 은 $+0.3174$, 앱은 $+0.4188$} 이다. 집안
최대치의 \textbf{5.6배 · 7.4배}. 판과 LODO 라는 이 실험실의 계기 둘이
\emph{내가 잰 것 중 제일 큰 전이 효과를 못 본다}.

\section{그래서 다른 집 밖 짝으로 검증했다}

집 밖에서 난 일을 집안 계기로 검증하려던 것이 틀렸다. 검증은 \textbf{다른
집 밖 짝}으로 해야 한다. 비게임 앱은 \texttt{entry\_friction} 이 100\% 이고
짝 안 부호가 $-0.566$ --- KR 의 $-0.320$ 보다 더 극단이다. 등록: 여기서도
가리면 무표지가 $+0.15$ 이상 오른다.

\textbf{통과했고 더 컸다.} $-0.1389 \to +0.2799$ ($+0.4188$). 짝 안 부호가
더 극단인 쪽이 더 크게 올랐다 --- 두 점이지만 방향이 맞다. 대조로 뺀
\texttt{goods\_scale}($+0.396$ · 부호 일관)은 $-0.1148$ 로 반대로 갔다.

\section{grp 이 산 것은 되돌림이었다}

두 짝 모두 축을 빼면 grp 이득이 $0$ 아래로 간다($-0.0263$ · $-0.0405$).
노트 413 이 grp 을 채택할 때 산 집 밖 값은 \textbf{새 정보가 아니라 이 축
하나가 낸 손해의 상쇄}였다. 노트 419 의 ``도메인 사이 분산 흡수'' 는
맞았고, 흡수하던 분산의 대부분이 축 하나에서 나왔다.

\section{배포 판정 --- 판이 원리상 못 보는 변경}

여기서 규약들이 서로 부딪혔고, 부딪힌 자리가 이 노트의 내용이다.

\textbf{먼저 전역 제거(학습에서도 빼기)를 쟀다. 기각이다.} 짝 12뽑기에서
판 $-0.0018$ · SD $0.0030$ · \textbf{양수 2/12} --- 노트 378 조건 ① 이
그대로 기각한다.

\textbf{그런데 집 밖 두 짝에서 실제로 잰 것은 전역 제거가 아니었다.} 재기
전에 짠 코드를 다시 읽어 보니 가림이 \textbf{예보 쪽에만} 걸려 있었다 ---
학습은 축을 그대로 쓴다. 그러니까 $+0.3174$ 와 $+0.4188$ 은
\emph{``안 배운 도메인에는 이 축을 안 쓴다''} 라는 좁은 규칙의 값이다.
이 규칙은 성질이 다르다.

\begin{center}
\begin{tabular}{lr}
\hline
계기 & 좁은 규칙의 값 \\
\hline
판 (11 도메인) & \textbf{정확히 $0.0000$} --- 원리상 못 본다 \\
LODO 열하나 & 문턱 넘은 게 셋뿐 · 노트 439 \textbf{잴 수 없음} \\
\quad 축 있는 일곱 평균 & $-0.0030$ \\
\quad 축 없는 넷 & $+0.0000$ (가드) \\
\textbf{집 밖 KR 만화} & $\mathbf{+0.3174}$ \\
\textbf{집 밖 비게임 앱} & $\mathbf{+0.4188}$ \\
\hline
\end{tabular}
\end{center}

판이 $0.0000$ 인 것은 잰 게 아니라 \textbf{증명}이다. 판이 채점하는
도메인은 전부 학습에 들어 있어 규칙의 가지가 아예 안 열린다. 코드로
확인했다 --- 집안 예보 최대 차 $0.00\mathrm{e}{+}00$, 집 밖에서는 열리고,
입력은 안 건드린다.

\textbf{채택한다.} 노트 431 의 LODO 다수결은 여기서 기각이 아니라
\emph{기권}이고(잴 수 없음), 노트 426 의 거울 --- \emph{전이를 확실히 사는
변경은 판을 조금 잃어도 넣는다} --- 의 조건이 다 찼다: 집 밖 짝
\textbf{둘이 같은 부호}, 하나는 사전 등록 뒤 \textbf{앞서 잰 것과 독립인
짝에서 복제}, 판 손해는 ``$2\times$짝SD 안''이 아니라 \textbf{$0$}.

\textbf{되돌림 조항}(노트 432). 되돌릴 상태의 이름은 \textbf{``예보 때도
\texttt{entry\_friction} 을 쓰는 F18''} 이다. 되돌릴 조건도 미리 적는다 ---
\textbf{다음에 새로 여는 집 밖 짝에서 이 가림의 효과가 $0$ 이하면
되돌린다.}

\section*{규약 --- 집 밖에서 난 일은 집 밖에서 검증한다}

이 노트가 규약 하나를 만든다. \textbf{집 밖 짝에서 나온 효과는 집안
계기(판 · LODO)로 검증하지 않는다. 다른 집 밖 짝으로 검증한다.} LODO
열하나가 \emph{같은 수집 · 같은 계기 · 같은 시기}에서 나오기 때문이다 ---
수집 단위로 눈금이 어긋나는 일은 집 안에서는 원리상 안 보인다. 노트 430 이
LODO 를 만들 때 짝을 셋에서 열하나로 늘린 것은 옳았지만, 늘어난 여덟이
\textbf{집 밖 짝의 대체물은 아니다}.

\section*{남은 것}

\begin{itemize}
\item \textbf{CN 이 왜 안 오르나.} 독을 안 먹었는데 $+0.2249$ 다. 남은
후보는 라벨 분포(SD $0.485$ · 왜도 $1.84$ 대 KR $0.619$ · $0.74$) 하나다.
\item \texttt{entry\_friction} 을 \textbf{빼는 대신 도메인별로 눈금을
다시 매기면} 집 안과 집 밖을 다 살릴 수 있나.
\item 나머지 넷도 짝 안 부호를 재서 \textbf{집 밖 전용 무해 축 집합}을
정할 수 있나 --- 앱에서는 \texttt{venue\_prominence} 가 $+0.010$ 으로 거의
비어 있다.
\item LODO 씨앗 셋$\to$아홉(문턱 $0.0108 \to \sim 0.006$) --- 네 노트째
미룸.
\end{itemize}

\end{document}
